\subsection{THE COMPLEMENTS' RELATION AND THE LEGENDRE-GAUSS MULTIPLICATION FORMULA}

One derives immediately from formula (2) of VII, p.~315, that, for every \(z \in C\) ,

\[
\frac {1}{\Gamma (z) \Gamma (- z)} = - z ^ {2} \prod_ {n = 1} ^ {\infty} \left(1 - \frac {z ^ {2}}{n ^ {2}}\right).
\]

Now the Euler expansion of \(\sin z\) (VI, p.~287, th. 2) shows that

\[
z \prod_ {n = 1} ^ {\infty} \left(1 - \frac {z ^ {2}}{n ^ {2}}\right) = \frac {1}{\pi} \sin \pi z;
\]

taking account of the functional equation (4) of VII, p.~315, one then sees that:

PROPOSITION 1. For every complex \(z\) one has

\[
\frac {1}{\Gamma (z) \Gamma (1 - z)} = \frac {1}{\pi} \sin \pi z\tag{8}
\]

(the complements' relation).

COROLLARY --- For every real t one has

\[
| \Gamma (i t) | = \sqrt {\frac {\pi}{t \sinh \pi t}} \qquad (t \neq 0)\tag{9}
\]

\[
\left| \Gamma (\frac {1}{2} + i t) \right| = \sqrt {\frac {\pi}{\cosh \pi t}}.\tag{10}
\]

Indeed, one deduces from (8) that \(\Gamma(it)\Gamma(-it) = \frac{i\pi}{t\sin\pi it} = \frac{\pi}{t\sinh\pi t}\) , and one has \(\Gamma(-it) = \overline{\Gamma(it)}\) ; similarly, (8) gives

\[
\Gamma \left(\frac {1}{2} + i t\right) \Gamma \left(\frac {1}{2} - i t\right) = \frac {\pi}{\sin \left(\frac {\pi}{2} + \pi i t\right)} = \frac {\pi}{\cos \pi i t} = \frac {\pi}{\cosh \pi t},
\]

and one has

\[
\Gamma \bigl (\frac {1}{2} - i t \bigr) = \overline {{\Gamma \bigl (\frac {1}{2} + i t \bigr)}}.
\]

Now let \(p\) be any integer \(> 0\) and consider the product

\[
f (z) = \Gamma \left(\frac {z + 1}{p}\right) \Gamma \left(\frac {z + 2}{p}\right) \dots \Gamma \left(\frac {z + p}{p}\right).
\]

By (3) (VII, p.~315), for every \(z \neq -n\) ( \(n \in \mathbf{N}\) ), \(f(z)\) is the limit of the product

\[
\begin{array}{c} \frac {n ^ {(z + 1) / p} n !}{\left(\frac {z + 1}{p}\right) \left(\frac {z + 1}{p} + 1\right) \cdots \left(\frac {z + 1}{p} + n\right)}. \\ \frac {n ^ {(z + 2) / p} n !}{\left(\frac {z + 2}{p}\right) \left(\frac {z + 2}{p} + 1\right) \cdots \left(\frac {z + 2}{p} + n\right)} \dots \\ \dots \frac {n ^ {(z + p) / p} n !}{\left(\frac {z + p}{p}\right) \left(\frac {z + p}{p} + 1\right) \cdots \left(\frac {z + p}{p} + n\right)} \\ = \frac {n ^ {z + (p + 1) / 2} p ^ {(n + 1) p} (n !) ^ {p}}{(z + 1) (z + 2) \ldots (z + (n + 1) p)} \end{array}
\]

and in particular \(f(0)\) is the limit of the product

\[
\frac {n ^ {(p + 1) / 2} p ^ {(n + 1) p} (n !) ^ {p}}{((n + 1) p) !}
\]

from which it follows that \(f(z) / f(0)\) is the limit of

\[
\begin{array}{c} \frac {n ^ {z} ((n + 1) p) !}{(z + 1) (z + 2) \dots (z + (n + 1) p)} \\ = z   p ^ {- z} \left(\frac {n}{n + 1}\right) ^ {z}. \frac {((n + 1) p) ^ {z} ((n + 1) p) !}{z (z + 1) (z + 2) \dots (z + (n + 1) p)} \end{array}
\]

which, by (3) (VII, p.~315), gives

\[
f (z) = f (0) z p ^ {- z} \Gamma (z).\tag{11}
\]

But one can write

\[
f (0) = \prod_ {k = 1} ^ {p - 1} \Gamma \left(\frac {k}{p}\right) = \prod_ {k = 1} ^ {p - 1} \Gamma \left(1 - \frac {k}{p}\right) = \sqrt {\prod_ {k = 1} ^ {p - 1} \Gamma \left(\frac {k}{p}\right) \Gamma \left(1 - \frac {k}{p}\right)}
\]

since \(f(0) > 0\) ; the complements' relation then gives

\[
f (0) = \sqrt {\pi^ {p - 1} / \prod_ {k = 1} ^ {p - 1} \sin \frac {k \pi}{p}}
\]

and since the product on the right-hand side is equal to \(p/2^{p-1}\) (VI, p.~284, cor. 1) one sees finally that:

PROPOSITION 2. For every complex number z not an integer \(\leqslant 0\) and for every integer p \textgreater{} 0 one has

\[
\Gamma \left(\frac {z}{p}\right) \Gamma \left(\frac {z + 1}{p}\right) \dots \Gamma \left(\frac {z + p - 1}{p}\right) = (2 \pi) ^ {(p - 1) / 2} p ^ {\frac {1}{2} - z} \Gamma (z)\tag{12}
\]

(Legendre-Gauss multiplication formula).

PROPOSITION 3. For every real x \textgreater{} 0 one has

\[
\int_ {x} ^ {x + 1} \log \Gamma (t) d t = x (\log x - 1) + \frac {1}{2} \log 2 \pi\tag{13}
\]

(Raabe's integral).

First we establish formula (13) for \(x = 0\) . Since \(\log \Gamma(x) \sim \log \frac{1}{x}\) as \(x\) tends to 0, the integral \(\int_0^1 \log \Gamma(x) dx\) converges. Further, the function \(\log \Gamma(x)\) decreases on {[}0, 1{]} (VII, p.~310); for every \(\alpha > 0\) one thus has

\[
\frac {1}{n} \sum_ {k = 1} ^ {q} \log \Gamma \left(\frac {k}{n}\right) \leqslant \int_ {0} ^ {\alpha} \log \Gamma (x) d x,
\]

\(q\) being the largest integer such that \(q / n \leqslant \alpha\) . Since \(\int_{0}^{\alpha} \log \Gamma(x) dx\) tends to 0 with \(\alpha\) and also \(\frac{1}{n} \sum_{k=q+1}^{n} \log \Gamma\left(\frac{k}{n}\right)\) tends to \(\int_{\alpha}^{1} \log \Gamma(x) dx\) as \(n\) tends to \(+\infty\) (II, p.~57, prop. 5) one has

\[
\int_ {0} ^ {1} \log \Gamma (x) d x = \lim _ {n \rightarrow \infty} \frac {1}{n} \sum_ {k = 1} ^ {n} \log \Gamma \left(\frac {k}{n}\right).
\]

But, by (12), the right-hand side of this formula is the limit of

\[
\frac {n - 1}{2 n} \log 2 \pi - \frac {1}{2} \frac {\log n}{n},
\]

whence

\[
\int_ {0} ^ {1} \log \Gamma (x) d x = \frac {1}{2} \log 2 \pi .\tag{14}
\]

Next we remark that from the identity

\[
\log \Gamma (x + 1) = \log \Gamma (x) + \log x
\]

one deduces, on integrating, that for x \textgreater{} 0

\[
\int_ {0} ^ {x} \log \Gamma (t + 1) d t = \int_ {0} ^ {x} \log \Gamma (t) d t + \int_ {0} ^ {x} \log t d t.
\]

But the integral on the left-hand side is also equal to \(\int_{1}^{x + 1}\log \Gamma (t)dt\) . Thus, by (14),

\[
\int_ {x} ^ {x + 1} \log \Gamma (t) d t = \int_ {0} ^ {x} \log t d t + \frac {1}{2} \log 2 \pi = x (\log x - 1) + \frac {1}{2} \log 2 \pi .
\]

